% !TeX root = ../../Book1.tex
\chapter*{序言}
\addcontentsline{toc}{chapter}{序言}

分析学，大抵总从极限开头。学习初等微积分时，我们往往先熟悉计算：怎样求极限，怎样求导，怎样积分。继续往前，才逐渐发现，计算背后还有许多问题需要回答。什么样的集合可以测量？函数列收敛到何处？极限能不能穿过积分号？局部的平均和微分信息，又怎样控制整体？到了偏微分方程、几何测度论或概率论，这些问题便常常成为论证的起点。

\ 

我编这套书，缘起于一种很具体的不满足。测度论、泛函分析、几何测度与 Sobolev 空间，往往分散在不同课程里；读者学会了各章的定理，却未必看清同一种办法为何反复出现。我希望把这些材料放在一条连续的路线上，让极限与紧性、局部化与尺度控制、变换与表示这几种方法，在不同问题中逐渐显出彼此的联系。这样安排，并不求百科全书式的齐全；我更关心一个结果从哪里来，后来又在什么地方派上用场。

\ 

全书分三卷。第一卷建立实分析的测度基础，并逐步进入几何测度、Sobolev 空间与有界变差函数理论。第二卷讨论 Fourier 分析、偏微分方程和调和分析，重点转向频率、振荡与算子的估计。第三卷从全纯函数出发，沿共形几何、势论、模形式与解析数论展开。具备相应基础的读者可以分别阅读三卷；第一卷提供的语言和工具，会在后两卷中经常用到。

\ 

本卷先从集合系、可测映射和测度延拓讲起，再由简单函数建立 Lebesgue 积分与各种收敛定理。随后引入所需的泛函分析工具，讨论 \(L^p\) 空间、Radon 测度及测度的弱收敛。中段把目光转向测度的局部结构：覆盖定理支撑测度微分，Hausdorff 测度、面积公式与余面积公式，则把小尺度上的观察与几何量的计算联系起来。最后一编用分布和弱导数处理不光滑函数，并进入 Sobolev 空间、有界变差函数与有限周长集合。前面的可测性、逼近和紧性工具，会在后半卷反复用到，只是所处理的对象和需要控制的量已有不同。

\ 

严格性在这里有实际用途。有限可加性与可数可加性不同，逐点收敛本身不足以保证交换极限和积分；测度不完备时，在零测集上任意修改函数，还可能破坏可测性。书中因此尽量把核心构造写到底：先核对良定义，再处理可数操作、无穷值和局部到整体的拼接；关键假设在哪里使用，也尽可能当场指出。例子和反例与证明同样重要，它们帮助读者辨认结论的边界，也说明那些看似繁琐的条件为何不能轻易略去。

\ 

本书面向学过单变量与多变量微积分、线性代数，并有基本证明训练的高年级本科生或研究生。读者应熟悉集合、映射、数列和函数列，能够使用 \(\varepsilon\)-\(\delta\) 语言；测度论与泛函分析不作先修要求。附录补充主线所需的度量空间、点集拓扑、多线性代数和泛函分析，遇到缺口时可以回查，但不代替这些学科的完整课程。关于高阶几何测度论与流的附录，则供读完相关正文后继续选读。

\ 

初次系统学习，宜按章节顺序阅读，尤其不要略过测度延拓、可测函数逼近和收敛定理的证明。已有测度论基础的读者，可以先看章首目标和章末习题，再决定哪些证明需要重读。带星号的单元供选读，基础附录可随正文需要查阅。章末习题均附解答或提示；建议先独立尝试，再用解答核对思路。未能做出一道题时，也不妨先记下卡住的位置，读过后面的内容再回来看。

\ 

读证明时，我建议保留三问：极限在哪个空间中取得？控制量从哪里来？哪条假设排除了失败的情形？若一个证明依赖截断、紧性或零测集处理，就把那一步单独写出来。读反例时则反过来，先找被删去的假设，再看原来的论证从哪一步失效，例子又怎样使结论不再成立。这样读过的定理，更容易在新的问题中认出来。

\ 

本书使用参考资料，一为核对定理的表述、假设和证明路线，一为指出继续阅读的入口。各编导读说明伴读书目的分工，章首和正文则在需要比较处理方式、说明出处或引导深入阅读时给出引用。进入几何测度与弱可微性以后，Lawrence Craig Evans、Ronald F. Gariepy 的《Measure Theory and Fine Properties of Functions》\cite{Evans2015Measure}可作为伴读与对照，帮助读者从另一种组织方式理解测度、可微性与函数的精细性质。书中的基础推导仍尽量自足，不以引文代替应当展开的论证。

\ 

参考资料不妨多备几种，但进入一个专题时，最好从相应章节的引文出发，选定一部专著持续读下去。熟悉其中一条论证路线以后，再比较别的讲法，才更容易看清不同假设、记号和证明安排各有什么用处。

\ 

本卷所讲的，仍只是继续学习的基础。概率论、算子理论、非线性偏微分方程和更高阶的几何测度论，都另有自己的问题与技术。我希望读者合上这本书时，除了掌握一些可用的结果，也逐渐习惯追问它们的条件，辨认它们的适用范围，并从局部的计算中看见整体的联系。

\ 

末了，录苏轼《题西林壁》赠读者：

\begin{poemquote}
横看成岭侧成峰，远近高低各不同。\\
不识庐山真面目，只缘身在此山中。
\end{poemquote}

\ 

\begin{flushright}
	R. EverStarine
	
	2026年9月4日
\end{flushright}
